\documentclass[11pt,letterpaper]{article}
\usepackage[T1]{fontenc}
\usepackage[utf8]{inputenc}
\usepackage{lmodern}
\usepackage[margin=1in]{geometry}
\usepackage{amsmath,amssymb,amsthm,mathtools}
\usepackage{booktabs,longtable,array,tabularx,enumitem,microtype,xcolor}
\usepackage{fancyhdr,listings}
\usepackage[colorlinks=true,linkcolor=blue!45!black,citecolor=blue!45!black,urlcolor=blue!55!black,breaklinks=true]{hyperref}
\hypersetup{pdftitle={Pairwise Certificates for Adaptive SAT Contraction},pdfauthor={Research prepared for Vladimir Reshetnikov},pdfsubject={Certified exponential-coefficient improvements, public formalization update, and structural barriers}}
\setlength{\headheight}{14pt}
\pagestyle{fancy}\fancyhf{}
\fancyhead[L]{\small Pairwise certificates for adaptive SAT contraction}
\fancyhead[R]{\small Research report}
\fancyfoot[C]{\thepage}
\setlength{\emergencystretch}{3em}
\setlist{itemsep=3pt,topsep=5pt}
\allowdisplaybreaks[2]
\lstset{basicstyle=\ttfamily\small,columns=fullflexible,breaklines=true,frame=single}
\newtheorem{theorem}{Theorem}[section]
\newtheorem{lemma}[theorem]{Lemma}
\newtheorem{proposition}[theorem]{Proposition}
\newtheorem{corollary}[theorem]{Corollary}
\theoremstyle{definition}\newtheorem{definition}[theorem]{Definition}
\theoremstyle{remark}\newtheorem{remark}[theorem]{Remark}
\newcommand{\poly}{\operatorname{poly}}
\newcommand{\pc}{\operatorname{popcount}}
\newcommand{\K}{\mathcal K}
\newcommand{\norm}[1]{\left\|#1\right\|_2}
\newcommand{\ceil}[1]{\left\lceil#1\right\rceil}
\newcommand{\floor}[1]{\left\lfloor#1\right\rfloor}
\newcommand{\pin}{\texttt{a7ae5ab55f96}}
% Exact values are independently checked by code/check_cached.py.
\newcommand{\CachedC}{29.81160894538347}
\newcommand{\SafeCachedC}{29.812}
\newcommand{\CachedBetaSquared}{\frac{15952347164548384057}{49000000000000000000}}
\newcommand{\CachedRadical}{\frac{15952347164548384057}{368725000000000000}}

\title{Pairwise Certificates for Adaptive SAT Contraction\\[5pt]
\large A smaller cached exponential coefficient, a low-degree polynomial-space improvement, and sharper obstructions to changing the exponent order}
\author{Research prepared for Vladimir Reshetnikov}
\date{October 26, 2026 (Pacific date)}
\begin{document}
\maketitle
\begin{abstract}
We review the twenty-nine refinement packets in ProveIt through \pin,
including the latest round-four placements, and incorporate the newly public
Anthropic formalization identified by the user. A closing check includes
the documentation update \texttt{7de9d09fecd1}. We obtain a deterministic
normalized Boolean tensor-contraction bound with time and space
\[
 (s+\mathcal L)^{O(1)}
 2^{(\CachedC\ldots+o(1))s/\sqrt{\lceil\log_2s\rceil}},
\]
improving the previous comparable cached time coefficient $30.891019$
and space coefficient $30.753913$ (rounded).
The new ingredient is a certificate for choosing among many contraction
plans: distribute the squared norm of an averaged operation support among
pairs of plans, then maximize each coupled pair exactly. This avoids the
exponential dependence on the number of plans in the preceding certificate.
It combines with the already available rectangular matrix backend and a
strict $10\to13\to7$ recursion that pays for every child entry.

For initial degree at most $\lceil s^{1/12}\rceil$, a separate streamed
construction gives polynomial-space coefficient $16.81223380188329\ldots$,
against $23.6661688071455\ldots$. We also extend the streamed flag
optimality obstruction to all binary separator trees, and prove a
graded-space rank obstruction suggested by the new formalization: a
polynomial-size crossing decomposition confined to energy $4h$ cannot
replace square-root graded entropy by a smaller positive power.

The new elementary certificates and resource calculations have proofs and
exact executable checks. Complete SAT consequences import the existing
contraction identities; neither a full solver nor a rebuilt Lean development
is supplied. The unrestricted SAT exponent order remains
$n\log\log n/\sqrt{\log n}$. Thus this is a strict improvement in explicit
exponential upper bounds, but not the requested improvement in the
logarithmic order of the worst-case exponent.
\end{abstract}
\newpage\setcounter{tocdepth}{1}\tableofcontents\newpage
\section{Outcome, comparison, and the frozen review}
\label{sec:review}

Let $s$ be the normalized wire budget, $\mathcal L$ the graph and
tensor-description length, and $\lambda=\lceil\log_2s\rceil$.
All numerical coefficients in this report multiply $s/\sqrt\lambda$ in a
base-two exponent. They are not coefficients per SAT variable.

\begin{center}\small
\begin{tabular}{p{7.7cm}rr}
\toprule Comparable setting & Prior coefficient & This report\\\midrule
General degree $\lceil s^{9/10}\rceil$, cached time &30.891019&29.811609\\
Same setting, cached space &30.753913&29.811609\\
Degree $\lceil s^{1/12}\rceil$, polynomial-space time &23.666169&16.812234\\
General degree, polynomial-space time &46.184846&46.184846 (retained)\\
\bottomrule
\end{tabular}\end{center}

The last row is unchanged. The new generic cached algorithm spends
subexponential space. The low-degree theorem has a stronger input premise
and must not be used as a generic $s^{9/10}$-degree result. Neither theorem
claims to dominate the smaller cached low-degree figures, which use more
space. The meaningful algorithmic envelope keeps every applicable older
branch and adds the two new ones.

If $U_{\rm old}$ and $U_{\rm new}$ are the displayed generic cached
upper-bound expressions, their ratio satisfies
\[
 \log_2(U_{\rm old}/U_{\rm new})
 =\left(30.891018920366\ldots-\CachedC\ldots+o(1)\right)
   \frac{s}{\sqrt\lambda}.
\]
This compares bounding functions. It is not a lower bound on the old
algorithm's running time or a measured SAT speedup.

\subsection{Latest relevant commits}

The GitHub connector supplied immutable source files and commit metadata.
The main review pin is \pin. The latest relevant history is:
\begin{center}\small
\begin{tabular}{p{3.0cm}p{10.7cm}}
\toprule Commit & Change\\\midrule
\texttt{2b925d8800cd} & Initial seven refinement packets placed in the new project.\\
\texttt{3e096089911f} & OpenAI preprint vendored from \texttt{7ef5c0cc5c}.\\
\texttt{6d3adc450615} & Eight second-round packets placed.\\
\texttt{450252e9d86c} & Five third-round packets placed; two replay scripts restored.\\
\texttt{b9941a7396e4} & Nine fourth-round archives committed in incoming.\\
\texttt{a7ae5ab55f96} & Those nine packets placed in Research; incoming reduced to its README.\\
\texttt{7de9d09fecd1} & Public Anthropic formalization linked; the research packets are unchanged.\\
\bottomrule
\end{tabular}\end{center}
The packet-placement commit is dated 2026-10-27 03:15:51 UTC, or October 26 at
8:15:51 PM Pacific daylight time. Its mathematical frontier is the
appropriate comparator; comparing only against a round-two or round-three
number would miss the current best claims \cite{index,placement}.
The closing documentation commit is dated 03:20:40 UTC (8:20:40 PM Pacific
on October 26). Its two README edits already correct the former statement
that the Anthropic development was unpublished \cite{closing}.

The review inventories all twenty-nine packet summaries and retrieves their
LaTeX, README, and proof-status sources. Detailed proof work concentrates on
the strongest bounds and their shared prerequisites. It is not a line-by-line
certification of every imported proof. The source ledger records 140 project
files and the additional primary sources used below; source blobs are
recorded without bundling redundant copies of the corpus.

\subsection{Complete packet inventory}

Here PS denotes polynomial space. Cached coefficients are time/space.
Unless qualified as low-degree, numerical coefficients use the generic
normalized degree premise. Some packets state more than the selected
contribution listed here.

\begin{longtable}{p{5.2cm}p{8.5cm}}
\toprule Packet & Relevant contribution or scope\\\midrule\endfirsthead
\toprule Packet & Relevant contribution or scope\\\midrule\endhead
\bottomrule\endfoot
\path{cauchy_mode_improvement} & Factorial Cauchy modes, slope three; same exponent order.\\
\path{eth_contraction_improvements} & Sharper taper/precision; PS $648.441$; streamed coefficient extraction.\\
\path{3sat_precision_refinement} & Precision $24MD/u\to6MD/u$; the $1/\log n$ denominator is conditional.\\
\path{sat_density_improvement} & Uniform density-sensitive exponent envelope.\\
\path{eth_density_refinement} & Independent intermediate-density improvement; same worst-case order.\\
\path{3sat_refinement} & Scalar parity lift and smaller degree cap; second contraction lacks a factorization interface.\\
\path{3sat_asymptotic_analysis} & Restricted barriers and missing degree/compression interfaces.\\
\path{eth_asymmetric_packet_balancing} & Asymmetric pair packing and creation-side degree tracking; PS $76.909$.\\
\path{3sat_round2_refinement} & PS $84.066$; actual-domain rank and Heisenberg/thinning barriers.\\
\path{eth_balanced_packing} & Near-unit padding and thin precision guard; PS $94.748$.\\
\path{eth_high_moment} & Independent $94.748$ certificate; degree-sensitive variants.\\
\path{eth_matching_refinement} & Independent $94.748$ certificate through matching-potential packing.\\
\path{eth_packet_refinement} & Improved pair loads and guards; PS $95.664$.\\
\path{eth_recursive_refinement} & Missing residue class in Cauchy modes; PS $133.207$.\\
\path{eth_joint_cauchy} & Full-orientation slope-four modes; weighted pair entropy and cached algebra.\\
\path{sat_adaptive_contraction} & Streamed separators, PS $55.709$; cached $36.750/34.024$.\\
\path{sat_frontier_improvements} & Variable arity, PS $63.857$; cached $38.484/34.275$.\\
\path{subexponential_sat_frontier} & PS $63.852$; cached $37.930/33.623$; exact kernel improvements.\\
\path{subexponential_sat_kernel_improvements} & Composed PS $66.605$ and cached $44.538/39.496$; extraction kernels.\\
\path{subsat_contraction_tree} & Paid binary table construction; cached $35.355/34.275$.\\
\path{sat_adaptive_bounds} & Current generic PS best $46.185$ from linear flags; cached $34.325/34.024$.\\
\path{sat_adaptive_tree_certificate} & Current generic cached best $30.891/30.754$ using three-plan support certificates.\\
\path{sat_adaptive_frontier} & Adaptive quartet PS $50.305$; cached $31.252/31.252$.\\
\path{sat_adaptive_matrix_partition} & PS $50.210$; cached $34.325/34.024$.\\
\path{subexponential_sat_adaptive_four} & Cached $32.916/31.662$; with a $19/8$ backend, $32.856/31.467$.\\
\path{subsat_rectangular} & Genuine rectangular backend, cached $33.994/33.994$; PS $53.978$.\\
\path{eth_adaptive_exact_kernels} & PS $55.332$, cached $35.071/34.000$; kernel-level timings.\\
\path{3sat_round3_frontier} & Best stated low-degree PS $23.666$; generic PS $63.757$.\\
\path{eth_adaptive_schedules} & PS $63.781$; cached $37.592/33.623$; incremental matching kernel.\\
\end{longtable}

The named round-three-frontier packet arrived with round four; it and the
adaptive-schedules packet had reviewed only through round two. The two primary
manuscripts are additional to these twenty-nine packets.

\subsection{The asymptotic target that remains open}

For fixed clause width $k$, the best proved order envelope reported in the
reviewed corpus is
\begin{equation}\label{eq:density-envelope}
 (n+m)^{O_k(1)}\,
 2^{O_k\left(\frac{n}{\sqrt{\log n}}
 \min\{\log(2+m/n),\log\log n\}\right)}.
\end{equation}
The direct clause-network bound can also be retained when it is smaller.
The general worst-case exponent is $O_k(n\log\log n/\sqrt{\log n})$,
and fixed density gives $O_{k,D}(n/\sqrt{\log n})$
\cite{source,density,density2}.

The conditional targets with a $\log n$ denominator remain conditional:
their missing structured summation or factorization algorithms are not supplied
by a smaller plan certificate. The separate Anthropic algorithm has the same
relevant logarithmic order and no stated comparable normalized-wire constant.
Consequently this report claims a strict improvement over the explicit
comparable cached certificates, not over an unknown big-$O$ constant or an
unproved algorithm.

\section{The inherited exact interface and new proof obligations}
\label{sec:interface}

\subsection{Input, rings, and immutable entries}

The input is a bipartite Boolean tensor network with at most $s$ wires,
polynomially many vertices, maximum degree at most $\lceil s^{9/10}\rceil$,
and description length $\mathcal L$. Its contraction is
\[
 Z=\sum_{x\in\{0,1\}^{E}}\prod_v T_v(x|_{\delta(v)}).
\]
Entries lie in $(\mathbb Z/2^{P_0}\mathbb Z)[V]/(V^\ell)$ with
$1\le P_0,\ell\le s$. A uniform entry routine takes
$(s+\mathcal L)^{O(1)}2^{O(s^{9/10})}$ time and polynomial space, with
fixed constants. The low-degree theorem changes the degree premise alone.

The source uses an extension ring
\[
 R_p=(\mathbb Z/2^p\mathbb Z)[U,V]/(f(U),V^\ell)
\]
with polynomial rank and precision. Each tensor is tagged with its creation
precision. Later evaluations use the same canonical integer representatives
of its creation-time value. Canonical lifting is not treated as a ring
homomorphism. These conventions, the deformation, and exact Hermite recovery
are imported \cite{source,asymmetric,joint,flags,trees}.

\subsection{Exact mode sum and entropy}

At an active call with $M$ wires, fixed alphabet $r$, retained taper $u$,
packet exponent $a$, and valuation $i<u$, put $D=2^{a-i}$. The compatible
inherited lemmas give child wire count $(2/r+o(1))M$, precision
$p_t\le(6+o(1))MD/u$, and at most $(2/r^2+o(1))M/D$ mode families of
each index on each directed pair. Their certified weights are
$a+1-i+4j-2\pc(j)$.

Enumerate an exact weight allocation $(b_{gh})_{g\ne h}$ with sum below
$p_t$. Let $\mathcal I_{gh}$ be its finite list of mode tuples. Skip an
allocation with an empty list and set
\[
 h_{gh}=\log_2|\mathcal I_{gh}|,\qquad
 \K(M,u)=\kappa M/\sqrt u,\qquad\kappa=\pi\sqrt2/\ln2.
\]
Its contribution is exactly
\begin{equation}\label{eq:exact-sum}
 \sum_{\xi_{gh}\in\mathcal I_{gh}}
       \prod_{g=1}^r Z_g((\xi_e)_{e\ni g}),
\end{equation}
where each $Z_g$ entry is a genuine recursive contraction. There are
polynomially many allocations since $r$ is fixed and all budgets are
polynomial in $s$. A knapsack recurrence computes exact list cardinalities,
rank/unrank, and polynomial-delay iteration; no counting oracle is introduced.

For every fixed nonnegative directed vector $z$, the common-tilt estimate is
\begin{equation}\label{eq:entropy}
 \sum_{g\ne h}z_{gh}h_{gh}
 \le \left(\frac{\norm z}{r}+o(1)\right)\K(M,u).
\end{equation}
To recall its leading constant, the mode weight generating product has
natural logarithm $\pi^2/(24\theta)+O(\log^2(2/\theta))$.
The shared budget then gives the following natural-logarithm bound,
before division by $\ln2$:
\[
 \theta p_t+\frac{\pi^2}{24\theta}
        \frac{(2+o(1))M}{r^2D}\sum z_{gh}^2.
\]
Optimizing the first two terms gives \eqref{eq:entropy}; the additive buffer
makes the error uniform on the fixed finite support family used here.

For unordered pairs $e=\{g,h\}$ write $g_e=h_{gh}+h_{hg}$.
Any unordered support $v$ therefore satisfies
\begin{equation}\label{eq:undirected-entropy}
 v\cdot g\le\left(\frac{\sqrt2\norm v}{r}+o(1)\right)\K.
\end{equation}
The factor $\sqrt2$ is essential. A fixed finite-net argument also gives
$\norm h\le(1+o(1))\K/r$ when needed below.

\subsection{Packing gates and the paid recurrence}

With $L=2^a$, $Q=r^a$, side degrees $R,C$, sufficient gates are
\begin{equation}\label{eq:gates}
 QR/M=o(1),\qquad Q^2C(\log s)^2/M=o(1),\qquad LQ^2/M=o(1).
\end{equation}
Initially both degrees are at most $s^{9/10+o(1)}$; afterwards the retained
side remains under that bound and the newly created side has degree at most
$2^{a_{d-1}+1}$. The low-degree proof tracks all previous packet degrees.

Each group has $2(r-1)$ incident directed pairs. Thus forming its full table
requires at most $2^{(\alpha_r+o(1))\K}$ recursive child evaluations, where
\[
 \alpha_r=\frac{\sqrt{2(r-1)}}r.
\]
If a certified local contraction plan costs $2^{(\beta_r+o(1))\K}$,
the correctly charged recurrence is
\begin{equation}\label{eq:recurrence}
 T_d\le\poly(s)\left[
    2^{(\alpha_{r_d}+o(1))\K_d}T_{d+1}
    +2^{(\beta_{r_d}+o(1))\K_d}\right].
\end{equation}
A smaller local operation coefficient alone is insufficient; the theorem
must control the maximum of all branches of this recurrence.

\subsection{Rectangular backend retained from the corpus}

We use the already reviewed field bounds $\alpha>3/10$ and $\omega<19/8$,
supported by the stronger primary results in \cite{matrix}. Fix
$\varepsilon=10^{-5}$. For sorted logarithmic matrix dimensions
$x\ge y\ge z\ge0$, the corpus's rectangular construction gives
\begin{equation}\label{eq:rectangle}
 F(x,y,z)=\max\left\{
 (1+\varepsilon/2)(x+y),\quad
 x+(47/56+\varepsilon)y+(15/28)z\right\}.
\end{equation}
The first term pays for input/output arrays. It cannot be omitted.
Tensoring a fixed $N\times N^{3/10}\times N$ algorithm of exponent
$2+\varepsilon$ with a fixed square algorithm of exponent $19/8$ gives
the second term, with square blocking of the largest dimension. Rational
coordinate lifting, nonexceptional odd primes, and CRT extend the algorithms
to $R_p$ without dividing by a nonunit. Fixed bilinear recursions process one
child product at a time and use a polynomial factor times the largest
input/output array in space \cite{rectangular}.

For completeness, the rectangle interpolation takes
$t_0=(y-z)/(1-3/10)$ and $t_1=(z-(3/10)y)/(1-3/10)$ when
$z\ge(3/10)y$. The arithmetic exponent, including blocking, is bounded by
$x-y+(2+\varepsilon)t_0+(19/8)t_1$.
When $z<(3/10)y$, padding gives $x+y+\varepsilon y$.
Both are at most \eqref{eq:rectangle}. The three matrix dimension
permutations are justified by the symmetries of the bilinear multiplication
tensor, with each permuted algorithm executed afresh.

The present contribution changes the adaptive plan certificate. It does not
claim the rectangular matrix bound, Cauchy basis, or Hermite recovery as new.

\section{The new weighted certificate for choosing a contraction plan}
\label{sec:portfolio}

\subsection{Every matrix operation has nine linear supports}

An operation merges disjoint sets $A,B$ with outside set $O$, or contracts
the last three blocks. Its three index classes are the unordered pair sets
$AB,AO,BO$. Their logarithmic dimensions are nonnegative numbers
$H_1,H_2,H_3$. Formula~\eqref{eq:rectangle}, after sorting these numbers,
is exactly the maximum of nine linear forms on the unsorted triple:
\[
 \operatorname{Perm}(1+\varepsilon/2,1+\varepsilon/2,0)
 \quad\text{and}\quad
 \operatorname{Perm}(1,47/56+\varepsilon,15/28).
\]
The first set has three vectors and the second six. This follows from the
rearrangement inequality, since $1\ge47/56+\varepsilon\ge15/28$.
They represent one sorted cost function, rather than nine separately
executed algorithms.

Multiply the supports by the fixed denominator $L_*=1{,}400{,}000$.
The integer triples become
\begin{equation}\label{eq:integer-triples}
 (1400007,1400007,0),\qquad(1400000,1175014,750000),
\end{equation}
with the permutations just specified. Assign each unordered group edge its
class coefficient, and assign zero to edges internal to a part.

A plan is an ordered partition of $[r]$ into three nonempty blocks. Build
each block by splitting its ordered list at $\lfloor k/2\rfloor$,
recursively building the two halves, then multiplying their tables.
Finally contract the three block tables. There are $r-3$ internal merges
and one final triangle, hence $r-2$ operations and $9(r-2)$ supports.
For plan $j$ let $\mathcal V_j$ be this finite set, and put
\[
 F_j(g)=L_*^{-1}\max_{v\in\mathcal V_j}v\cdot g.
\]
Up to polynomial factors, $2^{F_j(g)}$ pays for every local operation
after the singleton group tables have been computed.

\subsection{Couple the norm terms before taking pairwise maxima}

The preceding tree packet used all tuples of one support from each plan.
With $q$ plans and at most $m$ supports per plan, that exact certificate
has $m^q$ tuples. The following certificate needs only $\binom q2m^2$
support pairs. Its inequality can be weaker than complete tuple
enumeration, but its manageable cost permits a larger, more effective
menu of plans.

\begin{theorem}[Weighted pair certificate]\label{thm:pair}
Let $q\ge2$ and fix nonnegative integer weights $w_1,\ldots,w_q$ with
$W=\sum_iw_i>0$. For $i<j$ define
\begin{align}
 A_{ij}=\max_{\substack{v\in\mathcal V_i\\z\in\mathcal V_j}}
 \bigl(&w_i^2\norm v^2+w_j^2\norm z^2\notag\\
       &+2(q-1)w_iw_j\,v\cdot z\bigr),\label{eq:pair-max}
\end{align}
and
\[
 U=\frac{\sum_{i<j}A_{ij}}{(q-1)W^2},\qquad
 \beta^2=\frac{2U}{L_*^2r^2}.
\]
Selecting the plan of least $F_j(g)$ gives local nonrecursive time and
local storage at most a polynomial factor times
$2^{(\beta+o(1))\K(M,u)}$, uniformly over all exact mode-budget
allocations.
\end{theorem}

\begin{proof}
For arbitrary choices $v_i\in\mathcal V_i$, expansion of the squared norm
gives the exact identity
\begin{align*}
 \left\|\sum_i\frac{w_i}{W}v_i\right\|_2^2
 =\frac1{(q-1)W^2}\sum_{i<j}
 \left(w_i^2\norm{v_i}^2+w_j^2\norm{v_j}^2
       +2(q-1)w_iw_jv_i\cdot v_j\right).
\end{align*}
Each diagonal term occurs $q-1$ times. Taking the maximum of each pair
bracket separately can only increase the sum, so the squared norm is at most $U$.
For a given $g$, choose a maximizing support $v_i$ for each $F_i(g)$.
Then
\[
 \min_iF_i(g)\le\sum_i\frac{w_i}{W}F_i(g)
 =\frac1{L_*}\left(\sum_i\frac{w_i}{W}v_i\right)\cdot g
 \le(\beta+o(1))\K
\]
by~\eqref{eq:undirected-entropy}. The fixed finite menu makes the error
uniform. There is no exchange of a minimum and a maximum, and no assumption
that independently maximizing supports can be realized simultaneously.

Every table is an input or output of one of the plan's paid matrix
operations. Its logarithmic size is a sum of two dimension logarithms,
which is bounded by the first branch of~\eqref{eq:rectangle}. The
backend's space bound therefore covers all local tables, including the
singleton caches. Their costly recursive evaluation time is separate and
remains in~\eqref{eq:recurrence}.
\end{proof}

The weights are a proof device, not random choices at runtime. The
algorithm always selects the cheapest actual plan. Compared with separately
bounding every norm and inner product, \eqref{eq:pair-max} also keeps the
correlation between the norm and dot product of each pair of supports.

\subsection{The complete ten-group certificate}

Group labels start at zero. Each ordered block below is split as above.
The plan weights have denominator $W=1000$.

\begin{center}\small
\begin{tabular}{crlll}
\toprule Plan & $w_j$ & First block & Second block & Third block\\\midrule
1&137&$(7,4,3,9)$&$(1,6,0)$&$(8,2,5)$\\
2&124&$(5,4,6,7)$&$(9,1,8)$&$(2,3,0)$\\
3&135&$(0,5,9,7)$&$(2,1,4)$&$(8,3,6)$\\
4&120&$(0,4,7,8)$&$(2,9,6)$&$(1,5,3)$\\
5&108&$(3,1,7,2)$&$(9,6,8)$&$(5,4,0)$\\
6&119&$(7,1,5,8)$&$(2,9,4)$&$(0,6,3)$\\
7&106&$(7,2,8,6)$&$(5,9,3)$&$(0,1,4)$\\
8&87&$(1,9,6,7)$&$(2,5,0)$&$(4,3,8)$\\
9&64&$(4,6,1,5)$&$(8,9,2)$&$(0,7,3)$\\
\bottomrule
\end{tabular}\end{center}

There are 72 supports per plan, each in dimension $\binom{10}2=45$.
The exact checker reconstructs all 648 supports and evaluates
$\binom92\,72^2=186{,}624$ pair choices. It obtains
\begin{equation}\label{eq:beta10}
 \boxed{\displaystyle\beta_{10}^2=\CachedBetaSquared.}
\end{equation}
The JSON result contains every exact pair maximum and a maximizing support
pair, so the certificate does not depend on a floating-point search result.
The finite sum is reproduced from the plans and integer triples alone.

\subsection{Exact selection and complexity of the selector}

For each plan compute the integer
\[
 P_j=\max_{v\in\mathcal V_j}
 \prod_{\{g,h\}}
       \left(|\mathcal I_{gh}|\,|\mathcal I_{hg}|\right)^{v_{\{g,h\}}}.
\]
Minimizing $P_j$ is exactly equivalent to minimizing $F_j(g)$. Its bit
length is polynomial in $s$: $r$, $q$, the support coefficients and their
denominator are fixed constants, and the cardinalities have polynomial bit
length. Thus plan selection has no logarithmic comparison oracle or
rounding ambiguity. The literal integer-product routine in the code is
intended as an exact specification; the large fixed denominator makes it
unsuitable as a claim of practical efficiency.

The improvement is new relative to the reviewed corpus: a fixed finite
matrix-operation portfolio is certified by weighted coupled pairs, and
that certificate is integrated with a paid recursive schedule. Averaging,
matrix multiplication, and quadratic norm identities are classical; no
literature-wide priority claim is made for those ingredients.

\section{The complete cached time and space bound}
\label{sec:schedule}

\subsection{A strict rational schedule}

Put $\lambda=\lceil\log_2s\rceil$, $B=\lceil6\log_2\lambda\rceil$,
$u_d=\lfloor\eta_d\lambda\rfloor$, and
$a_d=\sum_{j=0}^d(u_j+B)$. Choose
\begin{align}
 r_0&=10,&\eta_0&=\frac{301}{20000},\notag\\
 r_1&=13,&\eta_1&=\frac{119}{10000},\notag\\
 r_d&=7\quad(d\ge2),&
 \eta_d&=\frac{1549}{324000}\left(\frac49\right)^{d-2}.
 \label{eq:schedule}
\end{align}
The cumulative leading fractions obey
\[
 \eta_0+\eta_1=\frac{539}{20000},\qquad
 \sum_{d\ge0}\eta_d=\frac8{225}.
\]
Stop at $H=\lceil s/\lambda^{3/4}\rceil$ wires and use the existing
packing construction with relative slack $\lambda^{-3}$. All active
tapers are positive eventually; small exceptional instances are handled
by exhaustive evaluation.

At depth one retain the old three-plan thirteen-group certificate
$\beta_{13}^2=1572/4225$ \cite{trees}. Its plans are
\begin{align*}
 &(0,1,2,3,4)\mid(5,6,7,8)\mid(9,10,11,12),\\
 &(12,10,1,6,4)\mid(9,0,3,8)\mid(5,2,7,11),\\
 &(0,8,11,12,1)\mid(7,10,4,3)\mid(5,2,6,9).
\end{align*}
These use the older $12/5$ square-block backend. In the tail retain the
single plan $(0,1,2)\mid(3,4)\mid(5,6)$ with
$\beta_7^2=632/1225$. The new rectangle portfolio is needed only at the root.

\subsection{Packing, depth, and precision are all paid}

The exact integer comparisons
\[
 10^{1000}<2^{3322},\quad13^{1000}<2^{3701},\quad7^{100}<2^{281}
\]
certify strict rational logarithm bounds. The dominant degree gates have
the following margins below one:
\begin{center}
\begin{tabular}{lll}
\toprule Stage & Exponent upper bound & Positive gap\\\midrule
Root &$9/10+2\eta_0(3322/1000)$ &$39/5000000$\\
Depth one &$9/10+(539/20000)(3701/1000)$ &$5161/20000000$\\
Tail &$9/10+(8/225)(281/100)$ &$1/11250$\\
\bottomrule
\end{tabular}\end{center}
The other fresh-side and packet-capacity gates are bounded by
$(8/225)(1+2(3701/1000))<1$. Thus every gate in~\eqref{eq:gates}
has a fixed positive exponent gap. It absorbs all polylogarithmic terms,
the guards, and $s/M\le\lambda^{3/4}$.

For a noncircular depth argument, first verify these gates provisionally
through $3\log_3\lambda/4+O(1)$ levels. In that range the smallest tail
taper is at least
\[
 \lambda^{\delta+o(1)},\qquad
 \delta=1-\frac{3\ln(9/4)}{4\ln3}>0.
\]
Nearly lossless packing then makes every tail child smaller than $M/3$
eventually. Hence an active path cannot exceed that provisional depth.
Floors are uniformly negligible and the relative padding errors multiply
to $1+o(1)$. In particular,
\begin{equation}\label{eq:masses}
 M_0\le s,\quad M_1\le(1+o(1))\frac s5,\quad
 M_d\le(1+o(1))\frac{2s}{65}\left(\frac27\right)^{d-2}\quad(d\ge2).
\end{equation}

Every retained valuation gives
$D\ge2^{a_{d-1}+B+1}$. The creation-time precision
$O(s\lambda^2 2^{a_{d-1}})$ is at most
$O(\lambda^{15/4}2^{-B})=o(1)$ times $MD/u_d$. Thus the sharp
precision and entropy estimates apply uniformly. All precisions have
polynomial bit length because $a_d\le(8/225)\lambda+O((\log\lambda)^2)$.
Original entry routines cost $2^{O(s^{9/10})}$ times a polynomial;
new packet routines cost $\exp(s^{8/225+o(1)})$; leaf work is
$2^{O(s/\lambda^{3/4})}$. Each is lower order on $s/\sqrt\lambda$.
The inherited immutable descriptions avoid recursively expanding entry
routines. Polynomial overhead over $O(\log\lambda)$ levels contributes
only $O(\lambda\log\lambda)$ to the exponent.

\subsection{Unroll the charged recurrence}

On the scale $s/\sqrt\lambda$, define
\[
 K_0^*=\frac\kappa{\sqrt{\eta_0}},\qquad
 K_1^*=\frac\kappa{5\sqrt{\eta_1}},\qquad
 K_2^*=\frac{2\kappa}{65\sqrt{\eta_2}}.
\]
The tail ratio is $\gamma=(2/7)/\sqrt{4/9}=3/7$ and
$\alpha_7/(1-\gamma)=\sqrt3/2>\beta_7$. Unrolling
\eqref{eq:recurrence}, including the terminal leaf route, gives the maximum
of three coefficients:
\begin{align}
 C_0&=\beta_{10}K_0^*
 =\frac\pi{\ln2}\sqrt{\CachedRadical}
 =\CachedC\ldots,\label{eq:C0}\\
 C_1&=\alpha_{10}K_0^*+\beta_{13}K_1^*
 =\frac\pi{\ln2}\left(\sqrt{\frac{7200}{301}}+
                       \sqrt{\frac{50304}{20111}}\right)
 =29.3352039993294\ldots,\label{eq:C1}\\
 C_\infty&=\alpha_{10}K_0^*+\alpha_{13}K_1^*
              +\frac{\alpha_7K_2^*}{1-\gamma}\notag\\
 &=\frac\pi{\ln2}\left(\sqrt{\frac{7200}{301}}+
          \sqrt{\frac{19200}{20111}}+
          \sqrt{\frac{77760}{261781}}\right)
 =29.06574852464125\ldots.\label{eq:Cinf}
\end{align}
Exact rational intervals in the checker prove $C_0>C_1,C_\infty$ and
$C_0<29.812$. Decimal output is not used in those comparisons.

Ancestor and descendant tables coexist on one active recursion path, so
their sizes add. Their exponent coefficients take a maximum, not a sum.
The local space coefficients are at most
$\beta_{10}K_0^*,\beta_{13}K_1^*,\beta_7K_2^*$, and the first is largest.
This proves a common time and space upper-bound coefficient. It does not
claim that the optimal space coefficient equals the optimal time coefficient.

\begin{theorem}[Generic cached bound]\label{thm:cached}
Under the explicit input premises and imported exact identities of
Section~\ref{sec:interface}, the construction is deterministic and exact,
with
\[
 T(s,\mathcal L),\ S(s,\mathcal L)
 \le (s+\mathcal L)^{O(1)}2^{(C_0+o(1))s/\sqrt\lambda}.
\]
In particular, $29.812$ is an eventual coefficient for each resource.
\end{theorem}
\begin{proof}
The chosen plan only reassociates the finite sums and products in
\eqref{eq:exact-sum}, so exactness follows from the imported identity.
The preceding gates justify its use at every depth. Theorem~\ref{thm:pair}
and the paid recurrence prove the time bound; the local table argument
proves the space bound. All remaining costs have just been included.
\end{proof}

\section{A low-degree polynomial-space improvement}\label{sec:lowdegree}

We import the adaptive linear-flag theorem of \cite{flags}: for a power-of-two alphabet $r$, its streamed local coefficient is
\[
 b_r=\frac1r\sqrt{\frac{(r-1)(4r+10)}9}.
\]
Its paid polynomial-space recurrence adds $(b_r+o(1))\K(M,u)$ at each network-recursion level.

For maximum initial degree $\lceil s^{1/12}\rceil$, the corpus's best
polynomial-space normalized coefficient is $23.6661688071455\ldots$ in
\texttt{3sat\_round3\_frontier} \cite{lowbaseline}. Composing adaptive linear flags with a
different compatible schedule yields
\[
 (s+\mathcal L)^{O(1)}
 2^{(16.81223380188329\ldots+o(1))s/\sqrt{\log_2s}}
\]
in polynomial space under the same imported contraction interface.
This is a coefficient improvement for the stronger low-degree premise,
not an improvement to the unrestricted logarithmic order.

Let $\lambda=\lceil\log_2s\rceil$, $g=\lceil\sqrt\lambda\rceil$,
$\Lambda=\lambda-g$, $b=\lceil6\log_2\lambda\rceil$, and $q=63/100$.
Choose alphabets and retained ranges
\begin{align*}
 (r_0,r_1,r_2)&=(32,16,8),&r_d&=4\quad(d\ge3),\\
 u_0&=\lfloor\Lambda/11\rfloor,&
 u_1&=\lfloor2\Lambda/99\rfloor,&
 u_2&=\lfloor2\Lambda/63\rfloor,\\
 u_d&=\left\lfloor(1-q)\frac{2\Lambda}{35}q^{d-3}\right\rfloor
 &&(d\ge3),\\
 a_d&=\sum_{j=0}^du_j+(d+1)b,& H&=\lceil s/\lambda\rceil.
\end{align*}
The cumulative leading retained fractions are $1/11,1/9,1/7$, and at
most $1/5$ in the tail. For each stage, $(1+2\log_2r)A=1$ at its indicated
cap $A$, so the negative $g$ margin controls the rounding gate $LQ^2/M$.
At the root the degree-packing exponent is
$10/11+1/12=131/132<1$. At the next two finite stages, using the immediately
preceding packet degree, it is $8/9+1/11=97/99<1$ and
$6/7+1/9=61/63<1$. In the tail it is at most $4/5+1/5=1$ and inherits
the negative $g$ margin. All old retained vertices have degree at most
the maximum of the original degree and earlier packet degrees; none gains
ports. Thus the same degree estimates cover both bipartite sides.

For hypothetical active depths at most
$\lceil\log_{20/11}\lambda\rceil+O(1)$, the minimum taper diverges because
$q>11/20$. Use relative padding slack $\lambda^{-3}$; the negative
$\sqrt\lambda$ margin absorbs its polylogarithmic cost.
All padding ratios are $o(1)$ uniformly. The child ratios are
$1/16,1/8,1/4$, then $1/2$, up to a multiplicative $1+o(1)$.
The tail ratio is eventually at most $11/20$, closing the depth bootstrap.
The total accumulated padding error is $1+o(1)$ on the $O(\log\lambda)$
levels. The precision buffer makes inherited precision lower order exactly
as in the common interface: $p/(MD/u)=O(\lambda^4 2^{-b})=o(1)$.
Uniformly diverging tapers absorb all floor errors.

The paid flag recurrence contributes $(b_r+o(1))K(M,u)$, where
$K(M,u)=(\pi\sqrt2/\ln2)M/\sqrt u$. Consequently the complete coefficient is
\[
 \frac{\pi\sqrt2}{\ln2}\left(
 \sqrt{11}\,b_{32}
 +\frac{\sqrt{99/2}}{16}b_{16}
 +\frac{\sqrt{63/2}}{128}b_8
 +\frac{\sqrt{35/2}}{512}
    \frac{b_4}{\sqrt{1-q}(1-1/(2\sqrt q))}
 \right)
 =16.81223380188\ldots.
\]
The exact rational checker proves this is below $16.813$.
Original entry costs, newly generated entry costs
$\exp(s^{1/5+o(1)})$, leaf costs $2^{O(s/\lambda)}$, polynomial arithmetic,
and the polynomial factors over the recursion are all lower order.
Creation precision remains immutable. The separator and network recursions
are streamed, preserving polynomial space.

The general-degree best $46.1848455867\ldots$ is not improved by this theorem.
The new coefficient is available when the original network satisfies the
additional degree premise, including bounded-degree direct CNF encodings.

\subsection{A concrete SAT application and its normalization}
The direct bounded-degree clause/equality encoding is an applicable case.
If $L$ is the number of literal occurrences and $d_v$ the occurrence count of
variable $v$, a bipartite construction with trivalent equality trees uses
\[
 M=L+2\sum_v(d_v-3)_+\le3L\le3km
\]
wires and degree at most $\max\{k,3\}$ \cite{flags}. Take
$s=\max\{M,n+1,2\}$ to pay for the exact-counting modulus $2^{n+1}$.
For fixed $k,D$ with $m\le Dn$, the new polynomial-space theorem therefore gives
\[
 (n+m)^{O_{k,D}(1)}
 2^{(16.81223380188329\ldots\,\max\{3kD,1\}+o(1))
                  n/\sqrt{\log_2n}}.
\]
This conservative explicit translation uses the existing encoding, not a
new wire reduction. It leaves the fixed-density logarithmic order unchanged.
For unrestricted density, the existing $s=O_k(n\log\log n)$ reduction
continues to have the same worst-case order after the new generic cached
contractor. Its unspecified wire constant prevents replacing $s$ by $n$
inside our numerical coefficient.

\section{Why more flexible separator trees do not improve the flag bound}\label{sec:treebarrier}

The round-four linear-flag packet \cite{flags} proved that
\[
 b_r=\frac1r\sqrt{\frac{(r-1)(4r+10)}9}
\]
is sharp among balanced binary separator trees under the Euclidean
pair-entropy relaxation, when $r$ is a power of two. The same obstruction
in fact applies to every binary separator tree, with arbitrary imbalance.

We work first with unordered pair weights. In a complete graph on $n$
vertices, let $F(S)$ denote the least possible maximum root-to-leaf sum of
cut weights in a binary separator tree on $S$. Exact finite distributivity
gives the recurrence
\[
 F(S)=\min_{A\sqcup B=S}\{w(A,B)+\max(F(A),F(B))\},\qquad
 F(\{v\})=0.
\]
Let $U(n)$ be this optimum when every edge has weight $1$. Let $A(n)$ be
the optimum when one distinguished vertex has incident edge weights $1$
and all other edges have weight $2/3$. By symmetry their recurrences are
\begin{align*}
 U(n)&=\min_{k+l=n}\{kl+\max(U(k),U(l))\},\\
 A(n)&=\min_{k+l=n}\left\{\frac{l(2k+1)}3+
                   \max\left(A(k),\frac23U(l)\right)\right\}.
\end{align*}
Every minimum above ranges over nonempty parts ($k,l\ge1$).
In the second line, the distinguished vertex lies on the side of size $k$.

\paragraph{Two auxiliary inequalities.}
For every positive integer $n$,
\[
 U(n)\ge\frac{n^2-1}3,\qquad
 A(n)\ge\frac{2n^2+2n-4}9.
\]
For the first, assume $l\ge k$ at a split and apply induction to $U(l)$;
the surplus over $(n^2-1)/3$ is $k(l-k)/3\ge0$.
For the second, put $G(n)=(2n^2+2n-4)/9$.
If $k\ge l$, use the $A(k)\ge G(k)$ branch. The surplus over $G(n)$,
multiplied by nine, is $l(2(k-l)+1)\ge0$.
If $k<l$, use $(2/3)U(l)\ge2(l^2-1)/9$ instead. The surplus is
\[
 2k(l-k-1)+l+2>0.
\]
The base case $n=1$ is exact.

\paragraph{Sharp value at powers of two.}
Write $L(n)=(2n^2+3n-5)/9$. We prove $A(2^j)=L(2^j)$ by induction.
For an unbalanced split of an even $n$, let $d=|k-l|\ge2$.
If $k>l$, the bound using $G(k)$ exceeds $L(n)$ by
\[
 \frac{2ld-k+1}9=\frac{d(2l-1)-l+1}9>0.
\]
If $k<l$, the bound using $U(l)$ exceeds $L(n)$ by
\[
 \frac{k(2d-3)+3}9>0.
\]
In the balanced case $k=l=2^{j-1}$, induction gives exactly $L(n)$;
the distinguished branch has the larger residual cost. The completely
balanced tree supplies a matching upper bound: every path has
$(n^2-1)/3$ unordered edges, and the path to the distinguished vertex
contains all $n-1$ heavier edges. Thus its weight is
\[
 (n-1)+\frac23\left(\frac{n^2-1}3-(n-1)\right)=L(n),
\]
and no other leaf path is heavier.

Return to directed pairs. Set $a_{gh}=1$ if $g=0$ or $h=0$ and $a_{gh}=2/3$
otherwise. Then
\[
 \|a\|_2^2=\frac{(r-1)(4r+10)}9=2L(r).
\]
Every binary tree has a path with $a$-weight at least $2L(r)$.
For $h=(K/(r\|a\|_2))a$, which lies on the Euclidean envelope
$\|h\|_2=K/r$, every binary separator evaluator therefore has worst path
cost at least $b_rK$. The imported adaptive flag upper bound attains it.

This rules out improving $b_r$ merely by allowing arbitrary unbalanced
binary separator trees at powers of two. It also covers separator nodes
of arbitrary arity: replacing any such node by a binary refinement causes
each refined root-to-child path to condition on a subset of the original
node's cross-child indices, so no path cost increases. It is a theorem about the stated
Euclidean entropy relaxation and child-evaluation tree model. It is not
a lower bound for the actual tensor instances, non-tree algorithms, or SAT.


\section{A rank obstruction for graded crossing decompositions}
\label{sec:graded-window}

This subsection concerns a concrete interface in Appendix B of the newly
public Anthropic development \cite{anthropic-paper}.  It is not a lower bound for SAT or for
arbitrary tensor algorithms.  All ranks are taken over a field $\mathbb F$.

For an alphabet $X$ of size $N$, write
\[
 \mathsf C(x,y,u,v)=[x=v][y=u].
\]
Its row flattening has rows $(x,y)$ and columns $(u,v)$; its column
flattening has rows $(x,u)$ and columns $(y,v)$.  Both flattenings are
permutation matrices of size $N^2$ and therefore have rank $N^2$.

\begin{proposition}[Finite graded-window obstruction]
Let $d\in\mathbb F$ be nonzero, and suppose
$d\mathsf C=\sum_{i\in I}A_i$.
Suppose that the row flattening of each $A_i$ factors through
$\bigoplus_{j\in J_i}W_{E_{ij}}$, where $W_E$ are finite-dimensional
graded spaces and $0\le E_{ij}\le Bh$.
Then
\[
 N^2\le \sum_{i\in I}\sum_{j\in J_i}\dim W_{E_{ij}}
 \le \Bigl(\sum_{i\in I}|J_i|\Bigr)
        \max_{0\le E\le Bh}\dim W_E.                 \tag{R1}
\]
The same assertion holds for the column flattening.
\end{proposition}

\begin{proof}
Factoring a matrix through a vector space of dimension $D$ gives rank at
most $D$.  Rank is subadditive under addition, and multiplication by the
nonzero scalar $d$ does not change rank.  Applying these statements to the
row flattening gives
\[
 N^2=\operatorname{rank}(d\mathsf C)
 \le\sum_i\operatorname{rank}(A_i)
 \le\sum_i\sum_j\dim W_{E_{ij}},
\]
followed by the maximum bound.  No commutation relation, positivity,
particular basis, or asymptotic partition estimate is used.
\end{proof}

\begin{proposition}[Modular version]
The finite rank inequality also holds if the pieces and their
factorizations are integral and reconstruction is required only modulo
$2^P$, provided $P>v_2(d)$:
\[
 \sum_i A_i\equiv d\mathsf C\pmod{2^P}.
\]
It is enough that each stated row factorization holds modulo $2^P$;
choose integer representatives of its factors and use their integer
products as the lifted pieces.
\end{proposition}
\begin{proof}
Let $B$ be the sum of those lifted row matrices, write
$d=2^v d'$ with $d'$ odd, and set $v=v_2(d)<P$.
Then
\[
 B=2^v\bigl(d'\mathsf C+2^{P-v}E\bigr)
\]
for an integer matrix $E$.  The matrix in parentheses reduces modulo
two to the invertible permutation matrix $\mathsf C$, so its determinant
is odd.  Thus $B$ has rank $N^2$ over $\mathbb Q$.
Rank subadditivity and the lifted factorizations give (R1), exactly as
before.  The column claim follows using the column factorizations.
\end{proof}
If $P\le v_2(d)$, the reconstructed right-hand side is zero and the
rank conclusion can fail, but that precision cannot recover even one
nontrivial bit of the crossing after removing the scalar $d$.


\begin{corollary}[The four-energy-window barrier]
Suppose $N=2^\lambda$, $h=h(\lambda)\to\infty$, and
$\lambda\sim c h^\alpha$, with $c,\alpha>0$.
Suppose the total number of summands in (R1) is polynomial in $h$ and
\[
 \log_2\!\left(\max_{0\le E\le4h}\dim W_E\right)
 \le c(4h)^\alpha+o(h^\alpha).
\]
Then $\alpha\ge1/2$.
\end{corollary}

\begin{proof}
Taking base-two logarithms in (R1) gives
$2\lambda\le O(\log h)+c4^\alpha h^\alpha+o(h^\alpha)$.
Divide by $c h^\alpha$ and take the limit: $2\le4^\alpha$.
\end{proof}

In particular, the hypotheses follow from regular power growth
$\log_2\dim W_E\sim cE^\alpha$ together with monotonicity of the dimensions,
and the input encoding $\lambda\sim\log_2\dim W_h$.
The stated uniform upper bound is the precise assumption actually used;
an asymptotic formula at the single energy $h$ alone is insufficient.

\subsection{Why paying for more pieces does not repair the same recursion}

There is a quantitative version of the obstruction.
Suppose each piece has at most $h^{O(1)}$ factorization channels,
$\lambda\sim c h^\alpha$, and the intermediate spaces still satisfy
the four-energy-window upper bound.  Then (R1) gives
\[
 \log_2|I|\ \ge\ (2-4^\alpha-o(1))\lambda.          \tag{R4}
\]
Thus a hypothetical $\alpha<1/2$ construction in the same window must
have exponentially many pieces, even if all identities are valid.

In the existing grid recursion, there are
$R=\binom g2^2=(1/4+o(1))g^4$ independently indexed crossing rectangles,
and the displayed implementation enumerates the Cartesian product
$I^R$.  On writing its normalized bit parameter as $N_{\rm grid}=g^4\lambda$,
the number of index tuples alone is at least
\[
 2^{\,((2-4^\alpha)/4-o(1))N_{\rm grid}}.
\]
This is a positive exponential floor when $\alpha<1/2$.
The conclusion is restricted to explicit independent enumeration of
piece indices.  A structured summation of those indices might evade it,
but it would be an additional algorithmic result; increasing the number
of pieces while keeping the current loop does not evade the barrier.


Appendix B uses $|I|=4h+1$ pieces, and each displayed factorization is a sum
of at most $2h+1$ intermediate spaces.  The intermediate energies are
$s+i$ and $2h-s-i$, with $-2h\le s\le2h$ and $0\le i\le2h$,
so every nonzero space has energy at most $4h$.
Thus the finite estimate applies with at most $(4h+1)(2h+1)$ summands.
For $0<\alpha\le1$, Appendix H's heuristic replacement
$\log\dim W_E\sim cE^\alpha$ would give the bit bound
\[
 c_i^{\rm row}+c_i^{\rm col}
 \le 2^{\alpha+1}\lambda+o(\lambda),
 \qquad \rho=2^{\alpha-1}.
\]
The corollary shows that simply transplanting $\alpha<1/2$ into these same
intermediate-energy and polynomial-piece interfaces is impossible.
This is stronger than the earlier corpus's obstruction \cite{roundtwo} that retained the
Heisenberg commutators: the present statement needs only a finite
factorization and rank subadditivity.  It does not assert that every
possible crossing decomposition has $\rho\ge1/\sqrt2$.

More generally, for a fixed $B>1$, the same rank proof forces $\alpha\ge\log_B2$ for an
intermediate-energy window $Bh$. A wider window changes the factorization
cost as well; it cannot be treated as a free relaxation.

\section{What the naive Cauchy--clique composition can certify}
\label{sec:nested-clique-obstruction}

Write $s$ for the normalized root-wire budget and $K\ge2$ for a growing
number of groups.  The existing packet layout uses $a$ coordinates,
retains $u\le a$ valuation levels, and has label-array size
$K^{2a}$.  Keeping this explicit layout polynomial in $s$ implies
\[
 u\log_2K\le a\log_2K=O(\log s).                    \tag{R2}
\]
The Cauchy mode enumeration certificate has logarithmic information
budget on the scale $M/\sqrt u$ at an $M$-wire call.
The newly public clique algorithm, as currently analyzed, has
$e(K)=\Theta((\log K)^{-1/2})$ on its permitted sequence of clique sizes.
Even if a compatible interface allowed this clique algorithm to multiply
the whole Cauchy information budget by $e(K)$ with no additional cost,
the resulting scale would be
\[
 \frac{M}{\sqrt u}\,e(K)
 =\Theta\!\left(\frac{M}{\sqrt{u\log K}}\right)
 =\Omega\!\left(\frac{M}{\sqrt{\log s}}\right).      \tag{R3}
\]
Here the last assertion is a statement about optimizing that proposed
analytic certificate, not a lower bound on actual runtimes.  It assumes
the displayed two-sided scales; an $O$-bound alone cannot prove a lower
bound.  It also grants the hypothetical composition an ideal interface:
in reality one must additionally pay for the local group-entry oracles,
their precision and integer representation, and the clique algorithm's
table-query convention.

Consequently the tempting nested invocation is not yet an order
improvement.  To make it one, at least one of the following quantitative
facts must change: the explicit layout cost in (R2), the Cauchy mode
information scale, the clique exponent as a function of $K$, or the
information actually presented to the clique engine.  Running the same
two estimates sequentially does not supply that change.


\section{The public Anthropic formalization and its scope}
\label{sec:formalization}

The supplied link identifies \texttt{anthropics/formal-math} PR~31 at
commit \texttt{1a6eee293b9e}.
GitHub metadata reports it merged on 2026-10-27 at 03:07:22 UTC,
or October 26 at 8:07:22 PM Pacific daylight time, with merge commit
\texttt{16a5ce41304f}
\cite{anthropic-pr}. ProveIt's closing commit \texttt{7de9d09fecd1}
already updates the earlier index's publication status and links this
development \cite{closing}. The new public paper has 900,337 bytes;
the older ProveIt copy has 868,673 bytes. They are distinct revisions.
The analysis of the graded crossing construction in this report uses the public
paper at the user's supplied commit.

We read the 49-page paper, \path{eth/README.md},
\path{eth/MainResults.lean}, \path{eth/Solution.lean}, the proof README,
the paper-to-Lean index, the comparator configuration, and toolchain pin.
\path{MainResults.lean} contains the endpoint definitions and imports
Mathlib. \path{Solution.lean} supplies theorems whose types are exactly
\texttt{MainResults.ETHIsFalse} and \texttt{MainResults.MainTheorem}
\cite{anthropic-main,anthropic-solution}.

The counting endpoint states one deterministic finite-alphabet multistack
Turing machine and one constant $C$, both chosen before $k$ and the input,
with time
\[
 2^{C\min\{km/\sqrt{\log_2m},\,
           k^2n\log_2\log_2n/\sqrt{\log_2n}\}}.
\]
Its domain is precisely specified: $k\ge1$, at least two clauses, at least
four variables, all declared variables occurring, and no repeated clause
lists. It counts assignments in a fixed bit encoding. The ETH endpoint
negates positivity of the deterministic 3-SAT exponent defined in that
machine model. These are substantive public statements, rather than merely
an unpublished proof claim.

The project pins Lean and Mathlib 4.33.1. Its README documents the build and
axiom-printing commands; the comparator configuration enables nanoda and
permits \texttt{propext}, \texttt{Quot.sound}, and
\texttt{Classical.choice}. These are source-reported verification procedures.
We did not run a fresh build, comparator, or kernel check of the transitive
development, and make no claim to have done so.

For fixed $k$, the published endpoint retains the same relevant logarithmic
exponent order. It does not supply a numerical normalized-wire coefficient
comparable to $C_0$. The proof README also explains that intermediate
counted operations may be polynomial-time functions of a bounded state,
with constants changing on conversion to actual machine steps. A numerical
constant from those intermediate counts therefore cannot simply be
multiplied into another contractor's bit-operation coefficient.

The link provides a different algebraic route worth exploring.
Appendix B represents crossing pieces using operators on a graded space;
their two factorization orders pass through complementary energies.
The four-energy rank theorem above uses that explicit structure to reject
one tempting but invalid route to a stronger denominator.

\section{Verification, deliverables, and logical status}
\label{sec:verification}

\subsection{What is proved here and what is imported}

\begin{longtable}{p{5.0cm}p{8.7cm}}
\toprule Item & Status\\\midrule\endfirsthead
\toprule Item & Status\\\midrule\endhead
\bottomrule\endfoot
Weighted norm-coupled pair certificate & New elementary proof; all finite maxima checked exactly.\\
Nine ten-group rectangle plans and weights & New fixed certificate; 648 supports and 186,624 pair choices regenerated.\\
$10\to13\to7$ time/space theorem & New compatible schedule and paid resource analysis, using stated imported identities and backend.\\
Low-degree coefficient $16.812234$ & New composition of the existing flag lemma with a verified degree-flow schedule.\\
All-binary-tree flag obstruction & New elementary proof; exact finite recurrences checked through size 128.\\
Four-energy-window rank obstruction & New finite rank proof and asymptotic corollaries; also valid for the specified modular reconstruction.\\
Cauchy deformation, exact recovery, packing, mode identity & Imported from pinned source and refinement packets.\\
Rectangular field algorithms and CRT bridge & Imported mathematical dependencies; not implemented by the finite checker.\\
Anthropic formalization & Public source and endpoints inspected; no fresh Lean build.\\
Complete SAT implementation or timing & Not supplied. No practical speedup measurement is claimed.\\
Smaller unrestricted logarithmic exponent order & Not established.\\
\end{longtable}

The complete contraction conclusions are conditional on the correctness of
the imported mathematical interface, in the same sense as the preceding
refinement packets. The new finite support inequality itself is unconditional
and applies to arbitrary finite vector menus.

\subsection{Exact computational checks}

All proof checkers use the Python standard library. They use explicit
exceptions for failed checks, so checks remain active under
\texttt{python -O}. The supporting results include:
\begin{itemize}
\item Exact regeneration of all nine rectangle support menus, the coupled
pair maxima and witness indices, and the rational value of $\beta_{10}^2$.
\item Exact root, depth-one, and tail packing inequalities, total taper
capacity, paid recurrence branches, and strict $C_0<29.812$.
\item Rational bounds on $\pi$, $\ln2$, and square roots: Machin's identity
with alternating arctangent bounds, the positive atanh series for $\ln2$
with a geometric tail bound, and integer square-root bounds. Decimal
renderings are for display only.
\item The streamed-tree lower recurrences and exact values at powers of
two through 128, together with the low-degree schedule and
$C_{\rm low}<16.813$.
\item Finite exact checks of crossing flattening ranks, modular full-rank
reconstruction, and negative controls showing why insufficient precision
cannot support that rank conclusion.
\end{itemize}

The new cached support construction was independently regenerated with a
separate implementation before the final certificate was selected. An
independent mathematical review checked the low-degree degree-flow proof
and the all-tree obstruction. The archive includes the final standalone
checkers, their outputs, and reproducible commands. Finite tests supplement
the symbolic proofs; they do not establish asymptotic uniformity by sampling.

\subsection{Reproduction and integration}

Run \texttt{python3 code/run\_checks.py} from the packet root to regenerate
the exact checks, then \texttt{bash build.sh} to compile the report with
pdfLaTeX. The main source is \texttt{report.tex}; all inputs are included.
No source repository or numerical optimizer is needed to replay the new
certificates. A fixed-data JSON file records the plans, weights, and
schedules, and the provenance directory identifies the reviewed commits
and source blobs.

The suggested integration destination is
\path{Computability/SubexponentialSAT/Research/sat_pairwise_certificates/}.
The integration note supplies a concise index entry and records that the
Anthropic publication update is already present at the closing head.
No GitHub mutation was made as part of
this review. The archive's integrity ledger excludes itself and temporary
LaTeX build files; it is not a claim that the imported corpus was formally
verified.

\section{Further research questions and the next structural targets}
\label{sec:research}

The current result reduces a coefficient, while the request for a smaller
order remains unresolved. The analysis gives sharper criteria for deciding
which next steps could actually change that outcome.

\paragraph{Actual image ranks of the crossing pieces.}
The four-energy obstruction rejects thinning the graded space while keeping
the same interfaces. It does not reject smaller factorizations through the
\emph{actual images} of the operators, provided a different, uniformly
proved factorization keeps the crossing identity intact. For a candidate
family $A_i$ on $\lambda$-bit strings, the target is
\[
 \max_i\left(\log_2\operatorname{rank}_{\rm row}A_i+
             \log_2\operatorname{rank}_{\rm col}A_i\right)
 \le4\rho\lambda+O(\log\lambda),\qquad \rho<1/\sqrt2,
\]
with polynomially many pieces, explicit integer factorizations,
polynomial-bit entries, and $2^{O(\lambda)}$ evaluation. An exact finite
rank test can reject bad candidates but cannot certify an asymptotic family.
The difficult new result would be a uniform cancellation or image-space
theorem, not a new Hilbert series substituted into the old proof.

\paragraph{More than two global regrouping directions.}
The Anthropic recurrence squares the clique size when it places the nodes
on a two-dimensional grid. A crossing identity compatible with three or
more useful global regroupings might change that recurrence. It would need
an explicit multi-direction identity, a balancing theorem, and counted
local evaluation. The two-order Heisenberg identity does not automatically
supply these. This is a structural research target, not a claimed algorithm.

\paragraph{Avoid explicit enumeration of correlated indices.}
Both the Cauchy mode sum and a hypothetical larger crossing-piece family
would benefit from structured summation of their indices. Merely encoding
an index by a short binary string does not give low-degree factors or a
cheap entry oracle. A proposed replacement must exhibit the factorization,
pay every coefficient query, retain creation precision, and prove a
strictly decreasing measure. The nested-clique calculation above explains
why composing the present two entropy estimates does not suffice.

\paragraph{Exploit correlations beyond the Euclidean envelope.}
The all-tree obstruction rules out improving the flag coefficient by
arbitrary tree imbalance alone. Actual generated mode factors could still
obey algebraic constraints absent from the relaxed entropy ball. A useful
lemma would show that the extremal star-heavy entropy profile cannot
arise together with costly child tensors, or would give an evaluator
sensitive to both remaining precision and mode weight. Simply lowering
requested precision leaves the inherited Hermite recovery increment and
does not by itself remove the square-root scale.

\paragraph{Improve certificates without confusing them with order gains.}
The pair certificate is a low-order relaxation of a finite support-choice
problem. Distributing the norm over triples or larger subsets, or choosing
another rational allocation of diagonal terms among pairs, can give tighter
certificates. The proof generalizes by counting how often each diagonal and
cross term appears. Such work could further lower explicit constants and
improve resource tradeoffs. As long as the mode entropy remains
$M/\sqrt u$ with $u=\Theta(\log s)$, it will not change the logarithmic order.

\paragraph{Formalize the small new lemmas first.}
The norm-distribution identity, finite support generation, all-tree recurrence
inequalities, and rank-window theorem are compact formalization targets.
They can be verified separately from the large analytic contractor. A later
end-to-end integration should then connect those lemmas to the exact ring,
packing, and machine-cost interfaces rather than treating a successful
finite checker as verification of the whole SAT algorithm.

\clearpage\begingroup\small
\begin{thebibliography}{99}

\bibitem{index} ProveIt. SubexponentialSAT project index and twenty-nine-packet frontier, frozen at the review pin. \href{https://github.com/VladimirReshetnikov/ProveIt/blob/a7ae5ab55f9627f5af7bf77443b3c34ef8ef1be4/Computability/SubexponentialSAT/README.md}{Pinned source}.

\bibitem{placement} ProveIt. Incoming placement record, including fourth-round arrival and relocation. \href{https://github.com/VladimirReshetnikov/ProveIt/blob/a7ae5ab55f9627f5af7bf77443b3c34ef8ef1be4/docs/incoming/README.md}{Pinned source}.

\bibitem{closing} ProveIt. Link the anthropics/formal-math Lean development of eth.pdf. Closing documentation-only commit \texttt{7de9d09fecd1}, October 27, 2026 UTC. \href{https://github.com/VladimirReshetnikov/ProveIt/commit/7de9d09fecd17de776d712835ecb2011205c531d}{Commit and changes}.

\bibitem{source} OpenAI. The Exponential Time Hypothesis is False, October 25, 2026. Quantitative reduction and contraction appendices; vendored from upstream commit 7ef5c0cc5c. \href{https://github.com/VladimirReshetnikov/ProveIt/blob/a7ae5ab55f9627f5af7bf77443b3c34ef8ef1be4/lib/openai-math/preprints/The-Exponential-Time-Hypothesis-is-False-October-25-2026/build/quantitative-contraction.tex}{Pinned source}.

\bibitem{asymmetric} Research packet. Asymmetric packet balancing: near-unit padding and fresh-side recursive degree bounds. \href{https://github.com/VladimirReshetnikov/ProveIt/blob/a7ae5ab55f9627f5af7bf77443b3c34ef8ef1be4/Computability/SubexponentialSAT/Research/eth_asymmetric_packet_balancing/report.tex}{Pinned source}.

\bibitem{joint} Research packet. Joint Cauchy modes: full-orientation integral separation and weighted pair-entropy estimates. \href{https://github.com/VladimirReshetnikov/ProveIt/blob/a7ae5ab55f9627f5af7bf77443b3c34ef8ef1be4/Computability/SubexponentialSAT/Research/eth_joint_cauchy/report.tex}{Pinned source}.

\bibitem{flags} Research packet. Adaptive bounds: linear-flag streamed summation, shared interface, and bounded-density wire count. \href{https://github.com/VladimirReshetnikov/ProveIt/blob/a7ae5ab55f9627f5af7bf77443b3c34ef8ef1be4/Computability/SubexponentialSAT/Research/sat_adaptive_bounds/report/report.tex}{Pinned source}.

\bibitem{trees} Research packet. Adaptive Contraction Trees with Exact Support Certificates: preceding cached frontier and the thirteen-/seven-group plans. \href{https://github.com/VladimirReshetnikov/ProveIt/blob/a7ae5ab55f9627f5af7bf77443b3c34ef8ef1be4/Computability/SubexponentialSAT/Research/sat_adaptive_tree_certificate/report.tex}{Pinned source}.

\bibitem{rectangular} Research packet. Rectangular multiplication with its input/output term and coordinate/CRT ring bridge. \href{https://github.com/VladimirReshetnikov/ProveIt/blob/a7ae5ab55f9627f5af7bf77443b3c34ef8ef1be4/Computability/SubexponentialSAT/Research/subsat_rectangular/sections/rectangular.tex}{Pinned source}.

\bibitem{lowbaseline} Research packet. Third-round frontier, received in round four: preceding low-degree polynomial-space coefficient. \href{https://github.com/VladimirReshetnikov/ProveIt/blob/a7ae5ab55f9627f5af7bf77443b3c34ef8ef1be4/Computability/SubexponentialSAT/Research/3sat_round3_frontier/report/frontier.tex}{Pinned source}.

\bibitem{density} Research packet. Density-aware exact SAT counting. \href{https://github.com/VladimirReshetnikov/ProveIt/blob/a7ae5ab55f9627f5af7bf77443b3c34ef8ef1be4/Computability/SubexponentialSAT/Research/sat_density_improvement/density_aware_sat_improvement.tex}{Pinned source}.

\bibitem{density2} Research packet. Density-sensitive refinement of the subexponential SAT bound. \href{https://github.com/VladimirReshetnikov/ProveIt/blob/a7ae5ab55f9627f5af7bf77443b3c34ef8ef1be4/Computability/SubexponentialSAT/Research/eth_density_refinement/report.tex}{Pinned source}.

\bibitem{roundtwo} Research packet. Round-two refinement: scoped Heisenberg and infinite-thinning barriers. \href{https://github.com/VladimirReshetnikov/ProveIt/blob/a7ae5ab55f9627f5af7bf77443b3c34ef8ef1be4/Computability/SubexponentialSAT/Research/3sat_round2_refinement/report/report.tex}{Pinned source}.

\bibitem{precision} Research packet. Sharper precision and the unproved structured-mode acceleration hypothesis. \href{https://github.com/VladimirReshetnikov/ProveIt/blob/a7ae5ab55f9627f5af7bf77443b3c34ef8ef1be4/Computability/SubexponentialSAT/Research/3sat_precision_refinement/report.tex}{Pinned source}.

\bibitem{anthropic-pr} Anthropic, formal-math pull request 31, at head commit \texttt{1a6eee293b9e}. \href{https://github.com/anthropics/formal-math/pull/31}{Pull request and merge record}.

\bibitem{anthropic-paper} Claude / Anthropic. A $2^{o(n)}$ Time Algorithm for 3-SAT, October 26, 2026. In particular Appendices A, B, D, E and H. \href{https://github.com/anthropics/formal-math/blob/1a6eee293b9e00cb432eef6018a21f523623a9af/eth/paper.pdf}{Paper at the supplied commit}.

\bibitem{anthropic-main} Anthropic formal-math. MainResults.lean: the independently readable formal endpoint definitions. \href{https://github.com/anthropics/formal-math/blob/1a6eee293b9e00cb432eef6018a21f523623a9af/eth/MainResults.lean}{Pinned definitions}.

\bibitem{anthropic-solution} Anthropic formal-math. Solution.lean and project verification documentation. \href{https://github.com/anthropics/formal-math/blob/1a6eee293b9e00cb432eef6018a21f523623a9af/eth/Solution.lean}{Pinned proof entry points}; \href{https://github.com/anthropics/formal-math/blob/1a6eee293b9e00cb432eef6018a21f523623a9af/eth/README.md}{README}; \href{https://github.com/anthropics/formal-math/blob/1a6eee293b9e00cb432eef6018a21f523623a9af/eth/proof/PaperIndex.lean}{Paper-to-proof index}.

\bibitem{matrix} Virginia Vassilevska Williams, Yinzhan Xu, Zixuan Xu, and Renfei Zhou. New Bounds for Matrix Multiplication: from Alpha to Omega. SODA 2024; arXiv:2307.07970v2, 2023. Its bounds $\omega\le2.371552$ and $\alpha\ge0.321334$ are stronger than the fixed inequalities used here. \url{https://arxiv.org/abs/2307.07970v2}.

\end{thebibliography}
\endgroup


\end{document}
